\gddchapter{Analysis}

\section{Thematic Analysis Framework}
Six phase model to move from theoretical raw audio to theoretical insights:
This is just a scaffold that will be filled out with actual data from this session:


% \begin{enumerate}
%     \item Familiarisation
    
%     Listen to the recorings while reading chapter 5 to notice intiial patterns

%     \item Initial Coding
    
%     Generate labels for specific UX moments such as "Command frustration", "modal confusion" etc.

%     \item Searching for Themes
    
%     Group codes into broader categories like "Neutral language friction" or "Immersion breaking"

%     \item Reviewing Themes
    
%     Check if those found themes accurately represent the difference between two input modalities, 
%     and if they fit the main research question.

%     \item Defining / Refining Themes
    
%     Refine the essence of what the data tells you about your research themes

%     \item Writing Up
    
%     Integrate transcript extracts in the finial presentation that show how the player 
%     interacted with the experience, including the themes defined above as anchoring points.

% \end{enumerate}


\subsection{Familiarisation}

Some initial patterns that were noted after listening to the recording:

% \begin{enumerate}
%     \item She registered the "Let's" just like \#2
%     \item The participant seems to be goal-oriented
%     \item She treated the descriptions very seriously, treating them as part of the puzzle not just navigation guidance
%     \item Klara asked twice if she can interact with the objects not listed in the description showcasing the curiosity and boundary-probing.
%     \item She attributed feeling overwhelmed to the novelty of the voice-driven experience rather than the modality itself.
%     \item Her answers on immerssion are inconsistent pointing to more nuance in her experience.
%     \item She framed the voice based navigation as the ChatGPT era and keyboard navigation as 80s Sierra text based games. 
%     \item The participant made remarks about seeking secrets hidden in the game and her desire to look for them.
%     \item She appreciated the writing style of the descriptions 
%     \item Klara was wearing a medical mask during the testing and expressed her concern over the accuracy of voice recognition.
 
    
% \end{enumerate}

\subsection{Initial coding}

Due to running out of time the initial coding was just me listening and reading the transcript without listing out those moments like in first 5 participants. 
         


\subsection{Searching for Themes}

\subsubsection{A - Foreign language friction}

The voice input modality introduced a to major sources of cognitive load. First one being the need to create and voice out input commands. The second cognitive load was the language of commands. For Maria English is a secondary language in which she's not as comfortable as she'd like to. That introduced a large pain point during the input stage. 

\subsubsection{B - Immersion vs usability}

Voice input produced a higher degree of immersion in the game at the cost of anxiety and higher cognitive load. This seems to be a trade-off observed with some previous participants.


\subsubsection{C - Open-ended input anxiety}

The freedom of open-ended voice input produced a feeling of anxiety in Maria, rather than feeling of control and agency.
The lack of clear feedback or UI scaffolding from the game made her feel even more confused in this situation.


\subsubsection{D - Keyboard modality comfort}

For Maria the keyboard version of the game was more comfortable not only due to immersion but also the fact that she had prior experiences with keyboard-driven games form her childhood. She remarked that this set of experiences made her feel more secure while playing the keyboard version of the game.


\subsection{Reviewing themes}


%     Check if those found themes accurately represent the difference between two input modalities, 
%     and if they fit the main research question.


Below for quick reference are the research questions from the thesis: 

\begin{enumerate}
    \item How would the participants describe the experience of using voice input in the presented
game environment?
    \item How do users feel about using the voice modality compared to traditional input?
\end{enumerate}

The themes are analysed to see if they fit the main research questions and if they accurately represent the difference between the two input modalities.


\subsubsection{A - Recognition anxiety as an immerison supressant}

% Szymon's constantily considered if the system would undrestand him in the first place. This constant anxiety effectively stopped him from fully immersing himself into the experience and made him more stressed compared to the keyboard version of the game.

This theme answers the research questions since it tackles the feelings of anxiety and immersion and the differences between those feelings in the both game versions.


\subsubsection{B - Tabletop voice model}

% The participant used the voice narration model that's typical in the tabletop roleplaying games. This was just how the researched assumed the players would be using the game. He the model and expectations for this type of voice interacion.
The way the player uses their voice crafts the relationship with the game and in turn their emotions and experience. This theme talkes directly about that so it's relevant to the research questions.

\subsubsection{C - Present state vs potential}

% The player expressed a very precise split between the present state of the game and it's potential. While he stated that in the current version the keyboard experience is more immersive he added that the voice-driven game would be more immersive for him if implemented according to his wishes.

This theme also ansers the research questions since it talks about the player experience in the voice version of the game, both in present state and a potential future state.

\subsubsection{D - Immersion through cognitive offloading}



As discussed in previous reports immersion is a very important feeling when talking about games. This answers the research questions since it's talking about the immersion directly and the relationship the player has between cognitive offloading and that feeling. 



\subsection{Refining Themes}

\subsubsection{A - Foreign language friction}

Playing the voice driven version introduced a major cognitive load for the participant stemming from two separate sources. 
Firstly, she noted that creating voice input commands was more cognitively demanding compared to selecting them from a list of options on a keyboard. In the keyboard version of the game all of available options were laid out to the player. The voice version didn't show them. The player had to create every command based only on the context cues. Maria admitted that she wouldn't preferred more guidance to make her feel less anxious and more at ease. 

The second source of heightened cognitive load was the fact that she had to articulate all of her commands in english. For the participant that language is not her primary choice. She prefers speaking in Polish or Italian, keeping English as an option that's used only when absolutely necessary. The linguistic barrier was something that plagued all of the other participants since none of them used english as their primary language. This on it's own is an interesting observation, since most of them do play games in english. It might appear obvious but reading and speaking are two different linguistic skills and it looks like the participants are much more comfortable reading than speaking in English.


\subsubsection{B - Immersion vs usability}

Voice input produced a higher degree of immersion in the game at the cost of anxiety and higher cognitive load. This seems to be a trade-off observed with some previous participants.

Despite the aforementioned increase in cognitive load and heightened anxiety the participant noted a higher perception while playing the voice-driven experience compared to the keyboard version of the game. 
It seems like in the sample group the theory that the player would always choose the more immersive experience doesn't hold ground. 

Maria voiced her desire for a hybrid system in which a voice-driven narrative is scaffolded by questions targeted at her and audio narration. Questions like "How do you want to proceed?", "Where do you wanna go from here?". It showcases her desire for having a clear sense of direction. 



This finding correlates with many previous participants.


\subsubsection{C - Open-ended input anxiety}

The freedom of open-ended voice input produced a feeling of anxiety in Maria, rather than feeling of control and agency.
She noted the feeling of not knowing if she's playing the game correctly.

This type of uncertainty if very interesting and something that wasn't considered during initial development. 
The anxiety got additionally heightened due to the fact that the game provided only basic feedback on failed commands. This led the participant to question her actions even further. A need for a better feedback system was already articulated by some participants before and this test only confirms that need further. This suggests that open-ended voice-driven game experiences would potentially benefit from precise feedback models and command scaffolding.


\subsubsection{D - Keyboard modality comfort}

For Maria the keyboard version of the game was more comfortable not only due to immersion but also the fact that she had prior experiences with keyboard-driven games form her childhood. She remarked that this set of experiences made her feel more secure while playing the keyboard version of the game. Those mental models from prior gaming experiences back when she was 13 clearly mapped to the text based game offering her the feeling of comfort and familiarity, contrary to the voice-driven experience which was novel and not as comfortable for her. 

This brings an important point of familiarity influencing the emotions of the participant. As noted before some participants noted that they felt emotions of excitement while playing the voice version not due to the game itself but rather the fact that it was a novel experience for them. Maria articulated it the most cleanly since she had direct experiences with text based games, that some of the participants lacked. 
These feelings of novelty and excitement will have to be looked upon and acknowledged in the main thesis document.







\subsection{Writing up}

\subsubsection{A - Narrative commanding}

From the jump Marcel used voice in a very poetic and creative way. His commands were long, detailed and narrative rich. They weren't just commands, rather a full blown narrative sketches. 

\begin{quote}
    "I'm going inside the clothing store but I'm trying to step inside the steps that I saw on the dusty floor so that I'm covering my tracks because fortunately we look like we have the same foot size so that's my game plan."
\end{quote}

These are the commands of someone treating the voice interaction modality as a creative medium, not just as a simple input tool. 
This interaction style was much closer to being an author much rather than just issuing commands to the computer. He often pushed for things thet were much beyond what's possible:

\begin{quote}
    "I'm entering the clothing store with small little hops like a bunny."
\end{quote}

The fact that the system could not always honour these elaborate requests left him feeling that the experience had something to be desired, but that didn't eliminate the pleasure he took from creating these commands (especially when they worked). The made him feel more present in the game in this way compared to traditional input. In the interview he stated that his requests were partially just his style of interaction and partially stress testing 

\begin{quote}
    "I was thinking about making an answer but in a fun way to somewhat test the system... Even if it didn't include or even mention nothing about my ambitions, it was way more fun in that sense. I was way more immersed in the story. I really cared about my friend."
\end{quote}




\subsubsection{B - Immagination boolean}



Marcel directly addressed that the voice version of the game required creativity from him in comparison to the keyboard version that could be "speed run", a much more robotic and predictable experience. He restated that the voice version forced him to be present, creative and concentrated on his experience. 

\begin{quote}
    "The voice-operated game left me with much more imagination in that sense that I was being creative about my actions — maybe just a little bit but still it made me feel more involved. And the second game was different in that sense that I've had clear actions that I could do and I could navigate the field much faster but I've had less use of my imagination in the second game. So I guess the main difference would be imagination in that sense."
\end{quote}

This goes hand in hand with other participants observations on how the voice interaction is heavier on the brain to sustain in this version of the experience due having to define the affordances available to the player. 

\begin{quote}
    "With the keyboard version I just was moving, just doing robotic tasks and basically speedrunning the game... I had to create the fun myself in trying to do the game faster. And that was my engagement."
\end{quote}


\subsubsection{C - Intent vs reponse gap}


Since the commands Marcel used were so nuanced and complicated the system couldn't fully comply with them the way he wanted. Even if the basic intent of a command was fulfilled the details that he wished couldn't be executed as this layer of precision was simply not implemented.

\begin{quote}
    "It gave me a sense of more freedom. But when I saw that my actions are fairly limited, I felt like my freedom was taken away from me a little bit... It was a little bit overwhelming because I was moving slow and processing a lot of stuff and thinking about what I'm going to say. But I think the freedom was winning in that battle."
\end{quote}

Similar to participant \#2 he through his commands he showcased that the game should have a way of semantic understanding of his commands, a way in which the player can define "how" does he want to do a said action. This in itself seems like a very difficult problem to solve, since at least in this example the boundaries of the system were kept ambiguous. This in turn can produce completely unpredictable input, not only unpredictably wide (since the player can ask any arbitrary command) but as demonstrated here arbitrarily deep as well. This is the first play test that showcased that other axis of chaos that get's introduced when a voice input modality is paired with a system that doesn't show it's boundaries. Players will probe it and want to do many many things, sometimes in a very particular way.




\subsubsection{D - Novelty as engagement fuel}

Marcel expressed that the voice input has been more engaging for him due to the fact that he had no prior experiences of using voice as a control input. He stated that it kept him more interested and more engaged with the game, pointing to the fact that he was used to conventional keyboard controls. 

\begin{quote}
    "I would choose the first one because of it being more interesting, more engaging, because I don't use normally any sort of voice input utilities and I don't play games that include that sort of thing. So it was something new for me, something more fun, more interesting, not the same old, same old."
\end{quote}

This on it's own is an interesting finding since it points to the fact that had he had more experience with voice input modalities he potentially wouldn't have been as engaged with the game. To put it more directly, the fact that he was engaged in the game more had nothing to do with the voice input's inherent nature and everything to do with it's novelty. 

\begin{quote}
    "If I played, let's say, a lot of games, then they can become kind of dull. So if something is new for me, or it's creative or interesting, I value that because I'm pretty worn out by some normal or usual type of games like the second one. It was way more similar to the games that I played in the past, I would say."
\end{quote}


This is something that was touched on by previous participants, who stated that the voice version was interesting for them despite it's clear issues partially due to that novelty factor. 
It makes it difficult to truly compare those two input modalities since the voice input modality tied to the game engine with AI is just not something that can be played in commercial games these days. Even people who use the voice input in games would partially feel the effects of novelty due to the fact that the game is built to understand natural language in comparison to simple pre-defined commands. 




\section{Language fluency consideration}

% Include the level of English language that the participant used and their percieved degree of fluency.
% Are they native speakers? What about their accent? Make sure to include that too.

The player is on a self proclaimed B2 level in English. It's her second language. Despite her level she doesn't feel comfortable using English in her daily life. 

\section{English Transcript}
GGabriel: OK, good.

Gabriel: In that case, let’s start recording.

Gabriel: And I’ll begin with a question.

Gabriel: How do you define yourself when it comes to video games? Do you consider yourself a gamer?

Gabriel: Do you not play at all, play a bit, play mobile games, or play a lot?

Maria: I don’t play at all.

Gabriel: You don’t play computer games at all?

Maria: But I used to play a lot when I was 13.

Gabriel: OK.

Gabriel: OK.

Gabriel: Come on, Casey, you can jump in.

Gabriel: Good.

Gabriel: Now, regarding voice assistants like Siri, Gemini, Alexa, or similar—what’s your attitude toward them?

Maria: I really like them and use them often.

Maria: Except Siri is the worst.

Maria: I have the worst experience with Siri, but I also like dictating text messages instead of typing them.

Gabriel: Siri is actually pretty terrible.

Gabriel: And do you have the new Siri on your phone, or is it not available yet in Poland?

Maria: I don’t know.

Gabriel: It’s like a TV remote with a button—like many colors on the edge, or not?

Gabriel: No, no, no.

Gabriel: Old TV.

Gabriel: Even the new one isn’t that good anyway.

Gabriel: And when you use a computer or phone, do you use any accessibility features—like screen readers, color filters, or similar?

Maria: I use a black-and-white filter, because I set up a shortcut so that when Instagram launches, the grayscale filter turns on.

Maria: It helps me avoid using Instagram as much as I used to.

Maria: That’s really good.

Maria: But it’s more about limiting my phone use, not a necessity.

Gabriel: Not a necessity, exactly.

Gabriel: So there’s nothing you use out of necessity?

Gabriel: No, no.

Gabriel: OK.

Gabriel: Let’s move on.

Gabriel: What I’m going to do now—we’ll do it here.

Gabriel: Here’s the country—look.

Gabriel: It’s a text game.

Gabriel: The idea is that it’s a post-apocalypse in Warsaw.

Gabriel: Your brother went to a shopping mall to get food and hasn’t returned—you think he might still be there.

Gabriel: The interface works like this: here is the name of the location you’re in, and the image is a rough visual representation of the place.

Gabriel: Here you see which floor you’re on—you’re currently on the ground floor.

Gabriel: Here is a long description.

Gabriel: The most important parts are the first sentence and the last sentence or last two sentences—they give hints about the location. The rest is more descriptive, poetic text.

Gabriel: If anything is unclear, you can ask me, because the language is quite poetic.

Gabriel: Coffee stopped brewing.

Gabriel: Oh, coffee.

Gabriel: Can I wet it?

Maria: Yes, I’m coming.

Let’s see how it looks.

Maria: Good morning, Mr. President.

Maria: Good morning.

Maria: Good morning.

Maria: Now I’m going to drink coffee.

Gabriel: And?

Gabriel: The text is quite small.

Gabriel: Sorry about that.

Gabriel: The difference in this game compared to a typical experience is that you don’t use a keyboard or controller—instead, if you want to do something, you hold the R key, speak, for example:

Gabriel: “Let’s go to the clothing store right now.”

Gabriel: Maybe not.

Gabriel: I don’t know why.

Gabriel: Or maybe it bugged the first time.

Gabriel: Anyway.

Gabriel: Sorry.

Gabriel: Just once.

Gabriel: OK.

Gabriel: And what if I want to go to a shop?

Gabriel: Because there’s some bug the first time.

Gabriel: OK.

Gabriel: And what if we go back?

Gabriel: And as you can see, you can use natural sentences—you don’t have to think in single words.

Maria: How should it be properly?

Gabriel: Yes, not just one word.

Gabriel: Correct.

Gabriel: OK.

Gabriel: As I said—the first sentence is an introduction, then a poetic description of the location, and at the end there are instructions about what you can do.

Gabriel: Just pass the money door.

Gabriel: Where you can go, what is connected to this place.

Gabriel: OK, and now you’ll play for about 10 minutes.

Gabriel: I’ll take some notes, but I’m not evaluating your performance.

Gabriel: That’s not the point.

Gabriel: I’m more noting when the system breaks.

Maria: Ah, OK.

Gabriel: OK, let me start a timer.

Gabriel: Click.

Gabriel: And one more thing I didn’t mention—you can use, create new items, and pick them up.

Gabriel: Your inventory is on the right side and items will appear there.

Gabriel: That’s everything from me.

Maria: I’m looking for my brother.

Maria: Yes?

Gabriel: That’s right.

Gabriel: He warned you before entering now.

Maria: I want to enter.

Maria: Yuck.

Maria: I think it’s a treasure.

Maria: I want to leave… OK.

Gabriel: The place written there is connected to two places:

Gabriel: Clothing store and small storage.

Maria: I want to go to the clothing store.

Gabriel: Everything OK?

Gabriel: Everything OK?

Maria: Learning helps, learning.

Gabriel: I didn’t say that.

Maria: This is a crowbar.

Maria: This is a crowbar.

Maria: I want to pick this crowbar.

Maria: Please take me there.

Maria: Not like that—I’m going, right?

Gabriel: Yes, but the doors are blocked.

Gabriel: And you need to say that the crowbar would open them.

Gabriel: But you don’t have the crowbar.

Maria: So…

Gabriel: I didn’t see the crowbar.

Gabriel: Please take me to the storage.

Gabriel: I want to leave.

Maria: I want to enter the small storage again.

Gabriel: I’m not sure… try again because it didn’t register.

Maria: Hello, I want to go to storage again.

Gabriel: The first sentence is the most important, and then maybe the last two sentences—the middle is poetic.

Maria: This is a metal lining.

Maria: This is first, second, third.

Maria: This is a head obstruction.

Maria: I want to take this crowbar to my shop.

Gabriel: Oh, here is text at the bottom.

Maria: Let’s go to the main entrance.

Maria: Let’s go to the shop.

Maria: We’re running.

Maria: I don’t know, maybe someone is chasing me, I have to run.

Gabriel: It’s taking a bit.

Gabriel: OK.

Maria: Get close to the door and try opening it using a key.

Gabriel: Oh.

Gabriel: The assistant didn’t understand.

Maria: Take me to the clothing store.

Maria: You can… go to the corridor.

Gabriel: I’m not sure the door opened.

Gabriel: Maybe you should say it without mentioning the door.

Gabriel: Yes, because I don’t know if the corridor is open.

Gabriel: Ah, OK.

Gabriel: You can just say where to go, because it seems like it suggests exiting.

Maria: Go back to clothing store.

Gabriel: Yes, now we can go to the corridor.

Gabriel: And now the system changes.

Gabriel: You’re no longer using voice.

Gabriel: Now you use keyboard keys.

Gabriel: One, two, three, four.

Gabriel: If other things appear there.

Gabriel: Just play for 10 minutes.

Gabriel: You can continue the game.

Maria: Good morning.

Let’s see what happens.

Gabriel: Yes, there’s an elevator there.

Gabriel: If you want to go further, you see the numbers?

Gabriel: Those are all the options where you can go.

Maria: I go up.

Gabriel: That’s just a description of the stairs.

Gabriel: If you press two, you go down—but it’s actually just stairs going up.

Maria: OK.

Gabriel: Now you’re on the first floor.

Gabriel: From here you can go elsewhere.

Maria: If I could just adjust this a bit longer.

Gabriel: You can’t go higher.

Gabriel: The third option is the lobby.

Gabriel: Crafting works by combining two items—you need to have prerequisites.

Gabriel: Maybe I’ll use matches.

Gabriel: You can see additional items highlighted.

Gabriel: Click.

Gabriel: Now you have matches and a stick.

Maria: Maybe I’ll craft something.

Gabriel: Yes, and there’s an arrow.

Gabriel: You select one item, press Enter, then another.

Maria: But that doesn’t work—how do you combine matches with a stick?

Gabriel: No, you can’t.

Gabriel: Exactly.

Maria: I think the right one is…

Maria: But I took the matches.

Gabriel: Yes.

Maria: Go to gun store, yes!

Maria: Immediately!

Maria: And I pick a pistol.

Gabriel: There’s a bug—you can pick up unlimited pistols.

Maria: Great.

Maria: OK, one right.

Maria: Go back to main.

Maria: Travel.

Maria: Go back to Korean snack store.

Gabriel: What is this?

Gabriel: Trouble.

Gabriel: Small corridor.

Gabriel: Trouble.

Maria: OK. Escape.

Maria: One.

Maria: One.

Maria: So one travel.

Maria: Fishing stall I think.

Maria: Why is it black?

Gabriel: Because there’s a fire in the bush.

Maria: OK.

Gabriel: It should be black.

Gabriel: OK.

Gabriel: We can end here.

Maria: Mhm.

Gabriel: It wasn’t about finishing the game within 20 minutes.

Maria: Mhm.

Gabriel: You were close anyway—good job.

Gabriel: Now I’ll ask a few questions about your experience.

Gabriel: One moment.

Gabriel: Now that you’ve played both versions, how do you feel?

Maria: The second one was better.

Maria: I had choices instead of speaking.

Maria: I preferred either full interaction by speaking like in a chat with AI, or using button commands.

Maria: I understood the game better with commands.

Gabriel: So the second version.

Gabriel: Was there anything surprising at the beginning?

Maria: I felt overwhelmed and didn’t know what to do.

Maria: I didn’t know how to start.

Maria: That was the hardest part.

Maria: But I also don’t have experience with such games.

Gabriel: That’s fine.

Gabriel: Don’t worry.

Gabriel: How would you describe the difference between the two games to someone who hasn’t played either?

Hmm.

Maria: The difference was that pressing keys felt more intuitive than speaking to the computer.

Maria: Speaking felt unnatural because I’m not used to telling a game what I want to do.

Gabriel: You mentioned it was overwhelming—can you explain more how using voice felt?

Maria: I felt a language barrier because I’m not as comfortable in English as I am with reading and pressing keys.

Maria: When I read text and commands, I fully understand them.

Maria: But speaking was much harder and more tiring.

Maria: If it were in Polish, it would be much more enjoyable.

Maria: If it were in a language I’m fully comfortable with, I would try again.

Gabriel: In which version did you feel more immersed?

Maria: The first one.

Maria: Because I didn’t have to read text—I just looked at the image and imagined the world.

Maria: I would prefer the first version if I didn’t have to read at all.

Maria: If I spoke, the computer could read the text to me.

Maria: Even in a mechanical voice.

Maria: That would be very valuable.

Maria: So it would read it to you?

Maria: Yes, it would read it to me.

Maria: And during the game, it would also guide me and ask questions like:

Maria: “What do you want to do now?”

Maria: It would guide me a bit.

Gabriel: OK, something like that.

Maria: Yes.

Maria: What are we doing now?

Maria: Then I would have choices and could respond.

Maria: I would also like to express emotional information—what I feel when entering a place.

Maria: Like “I feel afraid to be here” or “I have a bad feeling.”

Maria: It’s like in a movie when a character enters a room and senses something bad will happen.

Maria: Or something interesting.

Maria: I’d like more freedom of movement and expression.

Maria: Then I’d prefer that over selecting options, because options are limiting.

Maria: Even speaking would let me stay longer in a place and explore more.

Gabriel: How does the effort of forming spoken commands compare to pressing keys?

Maria: It’s definitely more effort with voice commands.

Maria: With keys it’s already there—I just press.

Gabriel: I understand.

Maria: That’s relaxing.

Maria: With voice, I have to think about what to say.

Maria: I don’t have enough data to feel comfortable.

Gabriel: Even though you had freedom to say anything, did it feel like more control?

Gabriel: Or was it more stressful?

Maria: It was more uncertain.

Maria: I didn’t know what I should do.

Maria: I felt lost.

Gabriel: And overall, how would you compare both experiences?

Gabriel: Which would you choose and why?

Maria: If nothing changed right now?

Gabriel: Yes.

Maria: I would choose the keyboard version because I understood it faster.

Gabriel: OK.

Maria: Because I had commands like “use item,” so I knew what to do.

Maria: I could go back and search rooms.

Gabriel: OK.

Maria: That was better for me.

Gabriel: OK.

Gabriel: And comparing the experiences overall?

Maria: The first version was more strange and unfamiliar, while the second was more familiar because I’ve used keyboard controls in games before.

Maria: It felt more intuitive.

Gabriel: OK.

Gabriel: At one point you asked if you can ask the game questions.

Maria: Yes.

Gabriel: What kind of questions?

Maria: What should I do now?

Maria: Or what’s behind these doors?

Maria: Something that expands the world.

Maria: Can I use a crowbar?

Maria: Is there another floor?

Maria: Questions that expand the map.

Gabriel: And when you used voice commands, what guided your choice of words?

Maria: Words I already knew.

Maria: I imagined keyboard commands—go forward, run, crouch, stay, go back.

Maria: Movement-based commands.

Gabriel: I understand.

Gabriel: And last question.

Gabriel: Is there anything from this research experience you’d like to add that we didn’t cover?

Maria: I would like more connections between rooms.

Maria: And I would like my own music.

Gabriel: Audio?

Maria: Yes.

Gabriel: Probably.

Gabriel: What do you mean by connections between rooms?

Maria: For example, earlier I couldn’t move from storage to another part of the mall.

Gabriel: So not going from A to C directly, but through B?

Maria: Yes.

Gabriel: OK.

Maria: I’d like more exploration like hidden paths and shortcuts.

Maria: I wanted to use the crowbar in unusual ways, like breaking a cabinet in the Korean store to discover something.

Maria: It made me feel a bit aggressive.

Maria: I wanted to interact more destructively.

Gabriel: I understand.

Maria: There were also no real threats.

Maria: Like maybe a dangerous dog in a room.

Gabriel: I understand.

Maria: Something I had to deal with.

Gabriel: Got it.

Maria: But I’m not sure I could handle that since I struggled already.

Gabriel: OK.

Gabriel: Thank you very much.

Gabriel: This was very valuable.


\section{Raw transcript}

Gabriel: Ok, dobrze.

Gabriel: W takim wypadku to niech się nagrywa.

Gabriel: I ja zacznę od takiego pytania.

Gabriel: Jak ty określasz siebie, jeżeli chodzi o gry komputerowe, czy określasz siebie jako gracza?

Gabriel: W ogóle nie grasz, albo grasz, albo mobilna też, ale czy grasz trochę, czy grasz bardzo?

Maria: W ogóle nie gram.

Gabriel: W ogóle nie grasz w gry komputerowej?

Maria: Ale dużo grałam, jak miałam 13 lat.

Gabriel: Okej.

Gabriel: Okej.

Gabriel: Chodź Casey, możesz skoczyć.

Gabriel: Dobrze.

Gabriel: A co się tyczy natomiast asystentów głosowych w postaci Siri, albo Gemini, albo Aleksy, albo czegoś takiego?

Gabriel: To czy jaki jest twój stosunek do nich?

Maria: Bardzo to lubię i często korzystam.

Maria: Tylko z Syrii jest najgorsze.

Maria: Z Syrii mam najgorsze doświadczenia, ale lubię też dyktować SMSy zamiast je pisać.

Gabriel: Syrii jest faktycznie straszna.

Gabriel: A ty w swoim telefonie masz taką nową Syrii, czy jeszcze w Polsce nie mamy jej?

Maria: Nie wiem.

Gabriel: To jest taka jak klikacz TV to jest taki jak guzik jakby takie jak wiele kolorów na krawędzi czy nie?

Gabriel: Nie, nie, nie.

Gabriel: Stara TV.

Gabriel: Znaczy nawet ta nowa nie jest jakaś dobra więc... I czy jak używasz komputera albo telefonu to czy używasz jakichś takich ułatwień dostępności na zasadzie czy ktoś czyta co jest na ekranie albo zmienia kolory albo tego typu rzeczy czy nie?

Maria: Używam filtra czarno-białego, bo zaprojektowałam sobie shortcut, że gdy uruchamia mi się Instagram, to włącza mi się ten filtr czarno-biały.

Maria: I to mi pomaga nie korzystać z Instagrama tak jak wcześniej tak długo.

Maria: Bardzo to jest dobre.

Maria: Ale to jest kwestia mojego ograniczenia korzystania z telefonu, a nie konieczności.

Gabriel: A nie konieczności, dokładnie.

Gabriel: Nie ma niczego, czego używasz przez konieczność.

Gabriel: Nie, nie, nie.

Gabriel: Dobrze.

Gabriel: No to w takim wypadku idąc dalej.

Gabriel: Teraz jakby to co ja zrobię, zrobimy to.

Gabriel: Tutaj jest kraj, o proszę.

Gabriel: To jest gra tekstowa.

Gabriel: Idea jest taka, że to jest post apokalipsa w Warszawie.

Gabriel: Twój brat wyszedł po jedzenie do galerii handlowej, do garnita mu nie wrócił i myśli, że tam dalej może być.

Gabriel: I interfejs wygląda w taki sposób, że tutaj jest nazwa miejsca, w którym jesteś, to obrazek jest jakby tego miejsca, taki poglądowy.

Gabriel: Tutaj jest na którym piętrze jesteś, na parterze aktualnie.

Gabriel: Tutaj jest opis, który jest dosyć długi.

Gabriel: Najważniejsze rzeczy w tym opisie to jest pierwsze zdanie i ostatnie tam zdanie, czy ostatnie dwa zdania, tam są jakieś wskazówki co do tego miejsca, powiedzmy reszta to jest taki opis lokalizacji, tak?

Gabriel: Jeżeli czego możesz mnie zapytać, jakby to było niezrozumiałe w tym, bo to jest taki język mocno poetycki.

Gabriel: Kawa przestała budkać.

Gabriel: O, kawa.

Gabriel: Przemoczyć?

Maria: Tak, już idę.

Zobaczmy, jak to wygląda.

Maria: Dzień dobry, panie prezydent.

Maria: Dzień dobry.

Maria: Dzień dobry.

Maria: To teraz będę piła kawę.

Gabriel: I co?

Gabriel: Tekst jest dosyć mały.

Gabriel: Przepraszam tutaj.

Gabriel: To co się różni, jako w tej grze typowego doświadczenia, to jest to, że nie sterujesz za pomocą jakiejś klawiatury albo kontrolera, tylko jeżeli chcesz coś zrobić, to klikasz tutaj dźwięk R, przytrzymujesz go, mówisz.

Gabriel: Na przykład, tutaj możemy pójść do miejsca, które się nazywa clothing store.

Gabriel: Napisane, że jest clothing store.

Gabriel: Trzeba mu powiedzieć, OK, how about we go to the clothing store right now.

Gabriel: Może nie.

Gabriel: Nie wiem dlaczego.

Gabriel: Albo za pierwszym razem jest bug.

Gabriel: Nieważne.

Gabriel: Sorry.

Gabriel: Tylko raz.

Gabriel: Okej.

Gabriel: A co jeśli chcę iść w sklep?

Gabriel: Bo za pierwszym razem jest jakiś błąd.

Gabriel: Okej.

Gabriel: Dobra, a co jeśli wracamy?

Gabriel: I jak widzisz możesz używać jakby, tak jakby, zdania naturałe, nie musisz jakby myśleć...

Maria: Jak to powinno być pogromnie.

Gabriel: Tak, nie ma tylko jednego słowa, tak?

Gabriel: Tak.

Gabriel: Pogromnego.

Gabriel: Okej.

Gabriel: No i okej.

Gabriel: Tak jak mówiłem.

Gabriel: Pierwsze zdanie jest jakaś taka wstępem, potem jest wzór takiego opisu, lokalizacji, takiego poetyckiego, a potem na końcu są zawsze informacje związane z tym...

Gabriel: Just pass the money door.

Gabriel: gdzie możesz pójść, co jest połączone z tym miejscem.

Gabriel: Okej, i teraz jakby przez 10 minut się może byś grała w to.

Gabriel: I ja sobie będę robił jakieś takie małe notatki, ale ja nie oceniam tego jak ci idzie, to nie o to chodzi zupełnie.

Gabriel: Okej.

Gabriel: Także tym się nie przejmuj, ja bardziej notuję rzeczy jak ten system coś się psuje w nim, na takiej zasadzie.

Maria: A, okej.

Gabriel: No, dobra, to czekaj, ja tu sobie odpalę timer.

Gabriel: Cyk.

Gabriel: A, i jeszcze jedno o czym nie wspomniałem.

Gabriel: używać i tworzyć nowe i podnosić.

Gabriel: I tutaj jest Twój ekwipunek po prawej stronie i one się tutaj będą pojawiać.

Gabriel: To wszystko z mojej strony.

Maria: Szukam brata.

Maria: Tak?

Gabriel: Tak jest.

Gabriel: Ustrzegł przed wejściem teraz.

Maria: I want to enter.

Maria: Fujka.

Maria: Obawiam się, że to skarb.

Maria: Chcę wyjść... O, tak, OK.

Gabriel: To miejsce, które jest napisane, to jest połączone z dwoma miejscami.

Gabriel: Clothing stores, klub z ubraniami i... Small storage.

Gabriel: Dokładnie.

Maria: I want to go to clothing store.

Gabriel: Wszystko okej?

Gabriel: Wszystko okej?

Gabriel: Mhm.

Maria: Nauka pomoże, nauka.

Gabriel: Nie stwierdziłem.

Maria: To jest crowbar.

Maria: To jest łom.

Maria: Mhm.

Maria: I want to pick this crowbar.

Maria: Please take me there.

Maria: Znaczy nie tak, tylko ja idę, prawda?

Maria: Mhm.

Maria: No to zrozumiem.

Gabriel: Tak, tylko że tutaj jest tak, że te drzwi są zablokowane.

Gabriel: Mhm.

Gabriel: I tego musisz mówić jakby, że łom by je otworzył.

Gabriel: Tylko, że tego łomu nie masz.

Mhm.

Maria: Więc...

Gabriel: Nie widziałam tu krowa.

Gabriel: Proszę odwiedź mnie do storażu.

Gabriel: Chcę wyjść.

Maria: I want to enter the small storage again.

Gabriel: Nie wiem czy... Spróbuj jeszcze raz się nagrać, bo chłopak się nie złapał.

Maria: Hello, I want to go to storage again.

Gabriel: Pierwsze zdanie jest, powiedzmy, to najważniejsze i później ostatnie dwa zdania może, a ten środek to tak bardziej poetycko jest.

Maria: To jest metalowe okładziny.

Maria: To jest pierwszy, drugi, trzeci.

Maria: To jest obiste gołowo.

Maria: Chcę zabrać tę krowę do mojej sklepówki.

Gabriel: O, tutaj jest taki tekst na dole

Maria: Idźmy do głównego wejścia.

Maria: Idźmy do sklepu.

Maria: Biegniemy.

Maria: Nie wiem, może ktoś mnie goni, muszę biec.

Maria: Nie biorę po małego.

Gabriel: Trochę zajmuję.

Gabriel: Ok.

Maria: Zbliż się do drzwi i spróbuj otworzyć je za pomocą klawiszu.

Gabriel: Oho.

Gabriel: Asystent się nie zrozumiał chyba.

Maria: Take me to the clothing store.

Maria: Możesz... Go to the corridor.

Gabriel: Nie wiem, czy on otworzył ten drzwi.

Gabriel: Może powinieneś... Możesz powiedzieć bez tego drzwi... Tak, bo nie rozumiem, czy korytarz nie jest otwarty aktualnie.

Gabriel: A, dobra!

Gabriel: Możesz powiedzieć bez tego, że gdzieś pójść, bo to chyba zasugerowało, żeby wyjść.

Gabriel: Jak mu powiesz, to jest bardziej o tym, żeby użyć tylko i wyłącznie.

Gabriel: Mhm, dobra.

Maria: Go back to clothing store.

Gabriel: Tak, czyli teraz możemy wejść na korytarz.

Gabriel: I teraz sprawowanie się zmienia.

Gabriel: Nie używasz już głosu.

Gabriel: Tylko używasz klawiszy na klawiaturze.

Gabriel: Jeden, dwa, trzy, cztery.

Gabriel: Jeżeli tam się jakieś inne rzeczy też pojawiają.

Gabriel: Przedstawia 10 minut po prostu.

Gabriel: Czyli możesz kontynuować grę.

Gabriel: Tylko już nie gadasz tylko sterujesz tak bardziej tradycyjnie.

Maria: Dzień dobry.

Zobaczmy, co się wydarzy.

Gabriel: Tak, tam gdzieś jest jakaś winda.

Gabriel: No i teraz jeśli chcesz stąd iść dalej, tak?

Gabriel: To tutaj widzisz jedynki?

Gabriel: Tutaj masz wszystkie opcje, w których możesz iść.

Gabriel: Idę na górę.

Gabriel: Tutaj jest taki mały... To jest tylko opis tych schodów, nie?

Gabriel: Ale tutaj...

Gabriel: Jak klikniesz dwa, to wejdziesz... To jest skarby going down, ale to są schody jakby wyżej po prostu.

Maria: Okej.

Gabriel: I teraz jesteś na pierwszym piętrze.

Gabriel: I stąd możesz już wyjść gdzieś indziej.

Maria: Jakbym sobie poprawiła tutaj to jeszcze dłużej.

Gabriel: Wyżej już nie możesz wejść niestety.

Gabriel: Dobra no to tam trzecia opcja to jest launch to po polsku jest właśnie takie lobby.

Gabriel: Crafting działa tak, że wybiera się dwa przedmioty.

Gabriel: Trzeba mieć załóżenie na tym, żeby to zrobić.

Gabriel: Może pójdę z tego karansnacka, bo tak to było wspomniane.

Gabriel: I tutaj widzisz, jak jest jakaś dodatkowa rzecz, to ona się też poświetla tutaj, czyli jakby... Widzisz, możesz podnieść się za pałki troszkę.

Gabriel: Tak.

Gabriel: Cyk.

Gabriel: Teraz możesz za pałki i łyk i punkt.

Gabriel: Super.

Maria: No to może kraków zrobię.

Gabriel: Tak, i tutaj jest taka mała strzałeczka wtedy.

Gabriel: Te strzałkami się wybiera, aby połączyć jeden przedmiot, Enter, z drugim.

Maria: Ale to nie działa, bo jak połączyć łące z zapałkami?

Gabriel: No nie, no właśnie, to się nie da.

Gabriel: No właśnie, no właśnie.

Gabriel: Dobra.

Maria: No to prawe jest to chyba.

Maria: Ale zabrałam zapałki.

Gabriel: Tak jest.

Maria: Go to gun store, tak jest!

Maria: Od razu!

Maria: I kapie pistol.

Maria: No i to mi się podoba.

Gabriel: To jest taki mały błąd, możesz podnieść tych pistoletów ile chcesz z podłogi i się nie rezentuje.

Maria: Świetne.

Maria: Dobra, jeden praweł.

Maria: Go back to main.

Maria: Travel.

Maria: Go back to Korean Snack Store.

Gabriel: Co to jest?

Gabriel: Find that much on the plate.

Gabriel: Find.

Gabriel: Paint.

Gabriel: Takie lekkie.

Gabriel: Rozmowne.

Gabriel: Tak.

Maria: Nie, ja tu nic nie zrobię, więc ja muszę tam wyjść i szukać brata.

Maria: Gdzie mam tu lunch?

Maria: Już tu byłam.

Maria: Traverse... Okej!

Maria: Dobra, ale stąd można iść, na przód, do... Jasrów Jednostek.

Maria: Mogę rozjaśnić go?

Maria: Dobra.

Maria: Pick up.

Maria: Legs stable right now.

Maria: Mhm.

Maria: Dobra, teraz wyjdźmy stąd.

Maria: I pójdźmy... do łazienki.

Maria: Ciekawe co tam jest.

Maria: Eee, suchy dno.

Maria: Intact?

Maria: Co to znaczy?

Gabriel: Tak, w całości.

Gabriel: Łazienka akurat tutaj nie ma zbyt wiele do odwiedzania łazience.

Gabriel: Łazienka jest pustym pomieszczeniem.

Maria: Dobra, to travel, return, go to lunch, travel, go to food supermarket.

Maria: O Jezu, a kto to?

Maria: Mój brat?

Gabriel: Tak, to on tam siedzi, no.

Maria: A, bo on poszedł po jedzenie.

Tak.

Maria: A, on skręcił nogę, więc ja mam legstabilizer dla niego.

Gabriel: Mhm.

Maria: No, to use item.

Gabriel: No i teraz powinniśmy strzokami wyglądać.

Gabriel: Aha.

Maria: Cyk.

Maria: Exit the building, okej.

Gabriel: Jeden.

Gabriel: Go to lounge.

Gabriel: I nie możesz... to jest błąd, że... ale to, że właśnie nie możesz wyjść w taki sam sposób, w jaki weszłaś.

Gabriel: Bo to się nie aktywuje.

Gabriel: Jest bug.

Gabriel: Ale domyślnie jest tak, że coś się dzieje i to domyślne wejście przestaje działać.

Gabriel: Musisz znaleźć inny sposób na wyjście z budynku.

Gabriel: Okej.

Gabriel: Dobra.

Maria: Trouble.

Maria: Konkurs podwóz.

Maria: No dobra, tu byliśmy.

Maria: To trochę znajomo wygląda.

Gabriel: Trouble.

Gabriel: Small Corridor.

Gabriel: Trouble.

Maria: A. Ok, Escape.

Maria: Jedynka.

Maria: Jedynka.

Maria: Więc jeden travel.

Maria: Fishing stall chyba to było.

Mhm.

Maria: Dlaczego jest czarne?

Gabriel: No bo jest płon w krzaku.

Maria: Okej.

Gabriel: Powinno być czarne.

Gabriel: No dobra.

Gabriel: Na tym etapie możemy sobie zakończyć.

Maria: Mhm.

Gabriel: To nie chodziło o to, żeby to koniecznie udało Ci się zakończyć grę w ciągu tych 20 dni.

Maria: Mhm.

Gabriel: No i tak byłeś blisko tego, że super.

Gabriel: Fajnie.

Gabriel: Dobra, teraz będę miał kilka pytań co do Twojego doświadczenia.

Gabriel: Poczekaj, tylko przeciągnę tutaj.

Gabriel: Dobra.

Gabriel: Teraz jak grałaś w obie wersje, to jak się czujesz?

Maria: Ta druga lepiej.

Maria: Miałam opcję do wyboru, a nie mówiłam.

Maria: Okej.

Maria: bo wolałam albo pełną interakcję i mówić jak ze sztuczną inteligencją w czacie, albo używać komend do wyboru.

Maria: Lepiej mogłam zrozumieć grę przy tych komendach.

Gabriel: Czyli przy tej drugiej wersji.

Gabriel: Czy było coś, co Cię jakoś zaskoczyło na początku?

Maria: Czułam się zakupiona i nie bardzo wiedziałam, co mam zrobić.

Maria: Jak zacząć?

Maria: Zacząć mi było najtrudniej.

Maria: Ale to też dlatego, że ja nie mam doświadczenia w takich grach.

Gabriel: Ale to spoko, to dobrze.

Gabriel: Nie przejmuj się.

Gabriel: Jakbyś miała opisać różnicę pomiędzy tymi dwiema grami?

Gabriel: Do osoby, która nie grała ani jedną, ani drugą.

Gabriel: To jakbyś opisała tą różnicę.

Hmm.

Maria: Różnica była taka, że bardziej intuicyjna była gra naciskając klawisze niż mówiąc do komputera.

Maria: Bo mówiąc intuicyjnie chcę użyć jakiejś komendy.

Maria: Nie jestem przyzwyczajona, że w grze mówię co chcę robić.

Maria: Jest to swobodna rozmowa.

Gabriel: Mówiłaś już trochę o tym, że to było przytłaczające, ale jakbyś mogła szerzej jeszcze opisać, jak to było używać głosu w grze?

Maria: odczuwałam barierę językową, bo nie czuję się tak swobodnie po angielsku, jak bym naciskała klawisze.

Maria: Jak mam napisany tekst i komenty, to to w pełni rozumiem.

Maria: A to mówienie było dla mnie dużo trudniejsze i dlatego się męczyłam bardziej.

Maria: Gdyby to było po polsku, to byłoby fajnie mówić.

Maria: Gdyby to było w języku, w którym się czuję zupełnie swobodnie, to chętnie bym jeszcze raz spróbowała.

Gabriel: A w której wersji czułaś się bardziej zanurzona w świecie przedstawionym?

Maria: W tej pierwszej.

Maria: Bo nie muszę czytać tego tekstu, tylko patrzę na ten obrazek i zaczynam sobie wyobrażać ten świat.

Maria: I wolałabym, żeby ta pierwsza wersja, żebym nie musiała w ogóle czytać tekstu.

Maria: Jeżeli mówię, to komputer da mi powiedzieć ten tekst, który jest napisany.

Maria: Nawet mechanicznym głosem.

Maria: To by było bardzo cenne.

Maria: Żeby ci przeczytał, tak?

Maria: Żeby mi to przeczytał, tak.

Maria: I żeby podczas tego grania, żeby ten głupio zadawał mi pytania.

Maria: Co chcesz teraz zrobić?

Maria: Żeby mnie trochę poprowadził.

Gabriel: Okej, okej, okej, no takie w zasadzie.

Maria: Tak.

Maria: To co robimy?

Maria: I wtedy mam wybór i mogę coś powiedzieć mu.

Maria: I chciałabym też udzielać takich informacji emocjonalnych, co czuję, jak wchodzę do jakiegoś wnętrza.

Maria: Boję się tutaj być.

Maria: Albo mam jakieś złe przeczucia.

Maria: Trochę jak w filmie, jak widzisz, jak bohater wchodzi do jakiegoś pomieszczenia i masz złe przeczucia, że coś tam się stanie.

Maria: Albo coś cię zainteresuje.

Maria: Albo co jest w tym...

Maria: pokaż mi, żeby była taka większa swoboda poruszania się po tym.

Maria: I wtedy wolałabym taką wersję niż wybieranie z opcjami, bo wybór opcji jest ograniczający.

Maria: Nawet wtedy takie granie z mówieniem, co chcę zrobić i jak się z tym czuję, dłużej bym była na jakiejś planszy, może bym więcej odkryła.

Gabriel: A jak wysiłek, który musiałaś wkładać w to, żeby sformułować te zdania, komendy, porównywalny jest z wysiłkiem, ok, który klawisz muszę kliknąć?

Maria: Tak, większy wysiłek jest w komendach, zdecydowanie.

Maria: Bo mam już napisaną opcję w klawiszu i wciskam klawisz.

Gabriel: Rozumiem.

Maria: I to jest takie relaksujące.

Maria: A budowanie komendy, zastanawianie się, co mam teraz zrobić.

Maria: Za mało mam danych, żeby swobodnie się w tym czuć.

Gabriel: To, że mogłaś powiedzieć cokolwiek, taka wolność niby do powiedzenia cokolwiek.

Gabriel: Czułaś, że miałeś większe poczucie kontroli?

Gabriel: Czy to było bardziej przejmujące niż...

Maria: Bardziej to było... powodowało większe... większą niepewność.

Maria: Co mam dalej zrobić?

Maria: Czy dobrze bram?

Maria: Zadałam sobie takie pytanie.

Maria: Nie widziałam od początku... celu.

Maria: Co mam...

Maria: Byłam zagubiona o tym.

Gabriel: A jak byś porównała ogólne doświadczenia w obu tych wersjach?

Gabriel: I którą z nich byś wybrała i dlaczego?

Maria: Gdybym nic nie zmieniała w nich, tak?

Gabriel: Tak, na teraz, no.

Maria: Na teraz wybrałabym grę klawiszową, bo szybciej zaczynałam rozumieć, o co chodzi.

Gabriel: Okej.

Maria: Bo mogłam, bo były komendy podane i jak była komenda użyj przedmiotu, to wiedziałam, że muszę coś znaleźć.

Maria: Mogłabym się cofnąć do jakiegoś pomieszczenia i poszukać czegoś, albo że coś pominęłam.

Gabriel: Tak.

Maria: To mi się bardziej podobało.

Gabriel: Ok.

Gabriel: A porównując doświadczenie pomiędzy jednym a drugim, tak?

Gabriel: Jak by się porównała to doświadczenie?

Gabriel: Rozumiem, że wybrałabyś tę wersję z tego powodu, ale samo doświadczenie?

Maria: Pierwsza wersja była bardziej dziwna, nieznana mi, a w tej drugiej wersji mogłam się odnaleźć, bo mam już doświadczenia w graniu gry, wciskanie klawiszy.

Maria: To jest dla mnie bardziej intuicyjny sposób porozumiewania się podczas grania.

Gabriel: Okej.

Gabriel: W którymś momencie na początku zadałaś mi pytanie też, jak grałaś, czy możesz zadawać pytanie grze.

Maria: Tak.

Gabriel: Jakie pytania chciałaś zadać?

Maria: Co mam teraz zrobić?

Maria: Albo takie, czy, nie wiem, czy ja chciałam zadać pytanie.

Maria: Co jest za tym, za tymi drzwiami?

Maria: Coś takiego, coś rozszerzającego to pole.

Maria: Czy mogę użyć łomu, tak, albo czy jest drugie piętro jeszcze?

Maria: Coś takiego rozszerzającego tą planszę, na której jestem.

Maria: Rozumiem.

Gabriel: A jak używałaś głosu do nominacji, to czym się kierowałaś przy wybraniu słów, których używałaś do mówienia?

Maria: Słowami, które znałam.

Maria: Jakbym wyobrażała sobie klawisze w grze, że idź do przodu, biegnij, padnij, zostań, wróć.

Maria: Takie ruchome.

Maria: Ruchowe komendy bardziej.

Maria: Czyli zamiast pisać komedy, to bym jej mówiła.

Gabriel: Rozumiem.

Gabriel: I to już ostatnie pytanie?

Gabriel: Czy jest cokolwiek dzisiaj, z naukowego doświadczenia, coś, co chciałabyś dodawać od siebie, czego nie poruszyliśmy?

Maria: Chciałabym, żeby było więcej połączeń między pokojami.

Mhm.

Maria: I żeby grała moja muzyka.

Gabriel: Jakieś audio?

Tak.

Gabriel: Pewnie.

Gabriel: Okej.

Gabriel: Przez połączenie między pokojami?

Maria: Co ty rozumiesz?

Maria: No chociażby to, co na początku, że nie mogłam, jak już weszłam do centrum, to nie mogłam z magazynu przejść dalej.

Gabriel: Że nie mogłaś przejść z punktu A do C, tylko musiałaś iść przez B, tak?

Maria: Tak.

Gabriel: Okej, okej.

Gabriel: Rozumiem.

Maria: I chciałabym, żeby było bardziej takich... i wtedy można by trochę poeksplorować to centrum, że znajdziesz jakąś... bo to mi przypominało doświadczenie w takich przygodowych grach, gdzie można by znaleźć jakieś skróty, jakiś komnat, przejście, jakieś ukryte skarby.

Maria: Chciałabym użyć łomu w jakimś nietypowym miejscu i na przykład miałam taką ochotę wziąć łom i stłuc to szafę w tym koreańskim sklepie i żeby może tam coś odkryć, spróbować.

Maria: Budziły się we mnie takie agresywne zachody.

Maria: Miałam ochotę zabić tego, zatrzymać tego w ramach.

Maria: Bardzo dobrze, bardzo dobrze.

Gabriel: No, no, no.

Gabriel: Rozumiem.

Maria: Brakowało mi też tam takich zagrożeń, że na przykład w jakiejś komnacie czeka na mnie pies chory na przykład.

Maria: No, no, no.

Maria: Rozumiem.

Maria: Muszę coś zrobić z tym.

Gabriel: No, rozumiem.

Maria: Ale nie wiem czy umiałabym sobie poradzić z tym, bo ciężko mi szło.

Gabriel: Spoko.

Gabriel: Dobrze.

Gabriel: Bardzo dziękuję za pytania.

Gabriel: To było bardzo cenne.

Gabriel: Dobra.







